\section{Related Work}
\label{sec:Related Work}

A detector that is accurate on its test set can still fail on the road, where it meets objects and conditions it was not trained for, and it keeps producing detections all the same. Such failures, which the safety of the intended functionality (SOTIF) addresses~\cite{iso21448}, call for a runtime monitor that tells the vehicle when to stop trusting its detector~\cite{rahman2021survey,yatbaz2024survey}. Monitors differ in the question they answer. Out-of-distribution (OOD) monitors ask whether the input contains an object of a class the detector was never trained on~\cite{du2022vos,du2022siren}. Uncertainty monitors and failure predictors ask how far the detector's output can be trusted~\cite{feng2022review,harakeh2021estimating,rahman2021perframe,yatbaz2024runtime}. Covariate-shift monitors~\cite{yang2024generalized} ask whether the input itself has changed while its content stays familiar, as when the image is foggy, blurred, noisy, or compressed~\cite{hendrycks2019benchmarking,michaelis2019benchmarking}. Our work belongs to this third kind. A monitor of this kind can take its signal from one of three places: the detector's predictions (Sec.~\ref{sec:rw-prediction}), the image (Sec.~\ref{sec:rw-image}), or the detector's features (Sec.~\ref{sec:rw-features}).

\subsection{Monitoring the Prediction}
\label{sec:rw-prediction}
Ideally, a model should be less confident on inputs it is not familiar with~\cite{hendrycks2017baseline}. The detector's outputs are therefore a natural place to start. SAOD~\cite{oksuz2023saod} averages the uncertainty of the three most confident detections to decide whether to accept an image. ContrastiveConf~\cite{park2026uncertainty} subtracts a weighted mean confidence of the queries a DETR~\cite{carion2020detr} discards from the mean confidence of those it keeps. Beyond the raw confidence, OOD scores computed from a classifier's logits, such as the energy~\cite{liu2020energy} and the maximum logit~\cite{hendrycks2022scaling}, follow the same idea. These scores react only once the corruption has disturbed the detections, so a mild corruption is easily missed (Sec.~\ref{sec:Experiment}).

\subsection{Monitoring the Image}
\label{sec:rw-image}
The image itself can be judged by a model beside the detector. Image-quality assessment (IQA) models predict how degraded an image looks. Image-based monitors differ in what they learn from: clean images, degraded images, or neither. Models that learn from clean images flag departures from them. NIQE~\cite{mittal2013niqe} measures how far an image's local contrast statistics differ from those of clean photographs. DisCoPatch~\cite{caetano2025discopatch} trains a discriminator to reject the imperfect patches of its own variational autoencoder, and flags an image whose patches the discriminator rejects. Models that learn from degraded images either recognize particular conditions from labeled examples, such as fog~\cite{pavlic2012image} or rain and snow~\cite{dhananjaya2021weather}, or learn, from synthetically distorted images and human ratings of them, one quality score that applies to any type of degradation, as BRISQUE~\cite{mittal2012brisque} and ARNIQA~\cite{agnolucci2024arniqa} do. CLIP-IQA~\cite{wang2023clipiqa} needs no quality-specific training: it rates an image zero-shot with a pretrained CLIP~\cite{radford2021clip}, by comparing it with the prompts ``Good photo'' and ``Bad photo'' All of these run a second model beside the detector and leave unused the features the detector has already computed.

\subsection{Monitoring the Detector's Features}
\label{sec:rw-features}
Finally, a monitor can take its signal from the detector's internal features. Hashemi et al.~\cite{hashemi2023runtime} fit a mean and spread to each neuron on clean images, and flag an image when more neurons than usual leave mean $\pm$ about two standard deviations. Becker et al.~\cite{becker2026operational} compare each channel's activation distribution in a frozen detector with the channel's distribution over the training set. A learned alternative fine-tunes a separate copy of the detector's backbone, with a small embedding head, on synthetically degraded images, and flags an image whose embedding lies far from that of clean images~\cite{becker2026degradation}.

% Not in your sketch of II-C; kept here until you decide:
% Geissler et al.~\cite{geissler2023low} classify injected noise, blur and contrast faults from a detector's channel sums with a decision tree.

Like Hashemi et al.~\cite{hashemi2023runtime} and Becker et al.~\cite{becker2026operational}, our monitor needs only the frozen detector and clean images. Their single global reference, however, mixes two sources of variation, the scene and the corruption. A mild corruption can thus hide within the differences between clean scenes. Our monitor therefore compares an image's early channels only with those of clean scenes of similar content, found by the detector's deepest stage, which reacts least to corruption yet still describes the scene (Sec.~\ref{sec:method}). It also separates two effects of corruption on these channels: a flattening, where their strongest responses fall, and a re-weighting, where their levels move up and down.
